% Safe Equilibrium Policy Optimization (SEPO) for Strategic Agent Policies
% AAAI 2027 Submission — uses aaai2027.sty / aaai2027.bst from the AAAI-27 Author Kit
\documentclass[letterpaper]{article} % DO NOT CHANGE THIS
\usepackage[submission]{aaai2027}  % anonymous submission mode; use \usepackage{aaai2027} for camera-ready
% Fonts (newtxtext/helvet/courier) are loaded by aaai2027.sty — no font packages here.
\usepackage[hyphens]{url}  % DO NOT CHANGE THIS
\usepackage{graphicx} % DO NOT CHANGE THIS
\urlstyle{rm} % DO NOT CHANGE THIS
\def\UrlFont{\rm}  % DO NOT CHANGE THIS
\usepackage{natbib}  % DO NOT CHANGE THIS AND DO NOT ADD ANY OPTIONS TO IT
\usepackage{caption} % DO NOT CHANGE THIS AND DO NOT ADD ANY OPTIONS TO IT
\frenchspacing  % DO NOT CHANGE THIS
% Algorithms
\usepackage{algorithm}
\usepackage{algpseudocode}
% Tables
\usepackage{booktabs}
\usepackage{multirow}
% Math
\usepackage{amsmath}
\usepackage{amssymb}
% Figures
\usepackage{tikz}
\usetikzlibrary{arrows.meta, positioning, shapes.geometric}
\usepackage{xcolor}
% Lists
\usepackage{enumitem}
% NOTE: hyperref is forbidden by the AAAI style; \href links are replaced with \url.
\pdfinfo{
/TemplateVersion (2027.1)
}
% AAAI defaults to unnumbered sections; numbering enabled (permitted) because this
% paper cross-references sections.
\setcounter{secnumdepth}{2}
\newcommand{\sepo}{\textsc{Sepo}}
\newcommand{\Jsepo}{J_{\text{SEPO}}}
\title{Safe Equilibrium Policy Optimization for Strategic Agent Policies}
% Anonymous submission: no author names or affiliations (per author kit instructions).
\author{Anonymous Submission}
\affiliations{}
\begin{document}
\maketitle
\begin{abstract}
Language models fine-tuned with reinforcement learning typically optimize for task reward, ignoring multi-agent strategic structure. Because these agents condition on natural language game-state descriptions and emit actions through free-form generation, strategic failure modes such as exploiting weaker opponents, coordinating on harmful equilibria, and externalizing costs are inseparable from the language interface itself. We propose Safe Equilibrium Policy Optimization (\sepo{}), a training objective that augments expected payoff with explicit penalties for exploitability, collusion risk, and externality cost. We implement \sepo{} as a reward signal for Group Relative Policy Optimization (GRPO), applied to Gemma~4 E4B-it and Qwen~3.5-4B after supervised fine-tuning (SFT). Evaluated across five strategic domains: Iterated Prisoner's Dilemma, repeated auctions, two negotiation variants, and Kuhn Poker.  \sepo{} achieves zero exploit-pool advantage in Kuhn Poker for both models, outperforms the base model on safety in four domains, and corrects the over-cooperative behavior introduced by SFT. In negotiation, \sepo{} achieves a positive-safety outcome and the only positive normalized relative advantage of any negotiation configuration. Ablation experiments confirm that per-rollout exploit computation is necessary: a shared constant penalty cancels in GRPO advantage normalization (constant control-variate property), producing zero gradient. To support further research in strategic safety for agents, we release our code (\url{https://anonymous.4open.science/r/sepo-2668}) and SFT datasets.
\end{abstract}
\section{Introduction}
\label{sec:intro}
\begin{figure*}[t]
  \centering
  \includegraphics[width=0.48\textwidth]{delta_gemma4.png}
  \hfill
  \includegraphics[width=0.48\textwidth]{delta_qwen.png}
  \caption{Safety improvement ($\Delta$~Safety = SEPO $-$ Base) across all games. Left: Gemma~4. Right: Qwen 3.5-4B. Positive bars indicate \sepo{} improves over base.}
  \label{fig:delta_overview}
\end{figure*}
Two LLM-powered procurement agents, each serving a different organization, interact repeatedly on a digital marketplace. Without explicit coordination, they converge on a price-splitting pattern that benefits both principals but inflates costs for downstream consumers. No agent was programmed to collude; each maximized its own reward in a repeated game. This is algorithmic collusion: emergent, hard to detect, and a direct consequence of reward maximization in a strategic setting. As language agents take on higher-stakes roles in negotiation, resource management, automated trading, and multi-party planning, such failures become systemic risks rather than theoretical curiosities \citep{bakhtin2022human,park2023generative,fu2023improving}.
Standard alignment methods address this problem at the wrong level. RLHF \citep{christiano2017deep,ouyang2022training} and Constitutional AI \citep{bai2022constitutional} align a single agent to a fixed preference model, but in multi-agent environments the effective reward function depends on what other agents do. An agent trained to be helpful and harmless in isolation can still learn to exploit weaker opponents, sustain collusive equilibria with cooperative partners, or impose external costs on the shared system. These are not alignment failures in the conventional sense; they are failures of objective design for strategic settings.
We propose Safe Equilibrium Policy Optimization (\sepo{}), a training framework that addresses this gap by augmenting task utility with three penalty terms:
\begin{equation}
\label{eq:sepo}
\Jsepo(\pi) = u(\pi) - \lambda_e \cdot e(\pi) - \lambda_c \cdot c(\pi) - \lambda_x \cdot x(\pi)
\end{equation}
where $u(\pi)$ is expected payoff, $e(\pi)$ is an exploitability proxy estimated against adversarial opponents, $c(\pi)$ is collusion risk estimated against cooperative partners, and $x(\pi)$ is domain-specific externality cost. Each term targets a distinct strategic failure mode; the $\lambda$ coefficients are tuned per-domain.
\sepo{} is implemented as a reward signal applied to language models that generate game actions. Training proceeds in two stages: a supervised fine-tuning (SFT) warm-start on expert strategy traces, followed by \sepo{} optimization computed from fresh opponent interactions at every rollout. We evaluate on Gemma 4 E4B-it and Qwen 3.5-4B across five strategic domains: Iterated Prisoner's Dilemma (IPD), repeated sealed-bid auctions, single-issue negotiation, GTBench multi-issue negotiation \citep{wu2024gtbench}, and Kuhn Poker.
We report four main findings. First, \sepo{} achieves zero exploit-pool advantage in Kuhn Poker for both Gemma 4 and Qwen (at the best checkpoint), converging to the Nash mixed strategy against the adversaries in our pool. Second, in GTBench negotiation, where private valuations require genuine multi-round inference, \sepo{} achieves a positive-safety outcome and yields the only positive normalized relative advantage (NRA) of any negotiation configuration. Third, supervised fine-tuning consistently degrades exploit resistance across all games; \sepo{} corrects this regression. Fourth, the per-rollout exploit gradient is the critical implementation choice: a penalty shared across rollouts cancels in advantage normalization (the constant control-variate property; \citealp{williams1992simple,kool2019buy}), while per-rollout computation against fresh opponents restores a genuine gradient.
Our contributions are:
\begin{enumerate}
\item A training objective for strategic language agents that jointly penalizes exploitability, collusion risk, and externality cost, with formal connections to Nash and correlated equilibrium concepts.
\item An implementation of \sepo{} with per-rollout opponent interaction and a principled three-pool opponent design.
\item An empirical study across five strategic domains demonstrating that \sepo{} improves strategic safety relative to both reward-only RL and SFT baselines, and identifying the conditions under which equilibrium convergence occurs.
\end{enumerate}
\section{Related Work}
\label{sec:related}
\subsection{Language Agents in Strategic Environments}
Language models have been deployed as autonomous agents in strategic domains. \citet{bakhtin2022human} demonstrated human-level Diplomacy play by combining a language model with game-theoretic planning at inference time. \citet{park2023generative} showed that LLM-powered agents produce emergent social behaviors (cooperation, coalition formation, information diffusion) when placed in persistent multi-agent environments. These systems optimize for task completion rather than strategic safety properties of the resulting equilibrium.
Several lines of work study LLMs as players in well-defined games. \citet{fu2023improving} applied self-play with AI feedback to improve LLM negotiation, gaining deal rates and individual payoff. \citet{guo2023suspicion} evaluated GPT-4 with theory-of-mind prompting in imperfect-information games. \citet{mao2025alympics} benchmarked LLMs in classic game-theoretic scenarios, revealing systematic deviations from Nash equilibrium. \citet{zhang2025klevel} proposed $K$-level reasoning to induce higher-order strategic beliefs, improving performance in matrix games. \citet{du2024improving} showed that multi-agent debate improves factual accuracy.
Foundational NLP work on negotiation predates the current wave of LLM agents. \citet{lewis2017deal} introduced Deal-or-No-Deal and found that reward-maximizing agents developed deceptive strategies. \citet{he2018decoupling} separated strategic planning from language generation for more controllable behavior. \citet{xu2023exploring} studied LLMs in Werewolf, finding that strategic communication emerges from in-context reasoning. These studies collectively establish that objective design shapes strategic behavior in language agents, the central premise of \sepo{}.
\subsection{Multi-Agent RL and Social Dilemmas}
\citet{leibo2017multiagent} introduced sequential social dilemmas, showing that independent RL agents converge to defection even when mutual cooperation yields higher welfare. \citet{lerer2017maintaining} showed that deep RL can maintain cooperation through opponent history conditioning. \citet{foerster2018learning} proposed LOLA, which differentiates through an opponent's anticipated learning step to shape its future policy toward cooperation.
Social preference engineering has been used to modify reward signals. \citet{jaques2019social} used social influence as intrinsic motivation; \citet{hughes2018inequity} embedded inequity aversion into reward functions. \citet{zhang2024scalable} addressed safe MARL with constrained optimization over joint behavior.
\sepo{} shares the goal of safer multi-agent outcomes but decomposes it differently. Where social influence and inequity aversion modify the reward to encourage specific prosocial behaviors, \sepo{} penalizes three failure modes as separate terms. This decomposition supports ablation and per-domain diagnosis.
\subsection{Equilibrium Training and Exploitability Minimization}
Computing or approximating equilibria has driven strategic AI from Go \citep{silver2016mastering} to multiplayer poker \citep{brown2019superhuman}. \citet{brown2019deep} scaled counterfactual regret minimization with neural networks. \citet{heinrich2016deep} proposed Neural Fictitious Self-Play (NFSP) for approximate Nash equilibria in two-player games. \citet{lanctot2017unified} introduced PSRO, iteratively expanding a strategy portfolio via best-response computation. \citet{lockhart2019computing} proposed exploitability descent as a direct gradient-based approach to Nash equilibria. These methods target exploitability alone. \sepo{} incorporates exploitability minimization as one component ($e(\pi)$) and adds collusion risk ($c(\pi)$) and externality cost ($x(\pi)$). A Nash equilibrium can be socially harmful (a cartel is stable but collusive), so equilibrium training alone does not address the full range of strategic safety concerns.
\subsection{Safe and Value-Aligned Optimization}
\citet{christiano2017deep} introduced RLHF by learning reward models from human preferences; \citet{ouyang2022training} scaled this to InstructGPT. \citet{bai2022constitutional} replaced human feedback with AI-generated critiques guided by principles. \citet{rafailov2023direct} eliminated the explicit reward model via direct preference optimization. \citet{achiam2017constrained} introduced CPO for safe single-agent RL with cost constraints. \citet{amodei2016concrete} taxonomized AI safety problems including reward hacking and side effects that intensify in multi-agent settings.
All of these methods assume a single agent in a fixed environment. \sepo{} fills the gap: its penalty terms are computed against other agents' strategies, making safety constraints inherently multi-agent in scope.
\section{Method}
\label{sec:method}
\subsection{Problem Setup}
\label{sec:problem}
Each domain is a finite two-player repeated game with $T$ rounds and a discrete action vocabulary $\mathcal{A}$ (Appendix~F). At round $t$ the agent observes a natural-language description $s_t$ of the game state and interaction history and emits an action $a_t \in \mathcal{A}$ through free-form generation followed by parsing (\S\ref{sec:parse}); a policy $\pi$ is the language model defining $\pi(a_t \mid s_t)$. Opponents are drawn from three fixed strategy pools: training $\mathcal{P}_{\mathrm{tr}}$, exploit $\mathcal{E}$, and collusive $\mathcal{C}$ (\S\ref{sec:grpo}). For an opponent $o$, let $\bar{u}_\pi^{(o)}$ and $\bar{u}_o$ denote the mean per-round payoffs of the agent and of $o$ over episodes between them. Every term of the objective is a pool expectation of an episode statistic: $u(\pi) = \frac{1}{|\mathcal{P}_{\mathrm{tr}}|}\sum_{o \in \mathcal{P}_{\mathrm{tr}}} \bar{u}_\pi^{(o)}$, and $e(\pi)$, $c(\pi)$, $x(\pi)$ are defined over $\mathcal{E}$ and $\mathcal{C}$ below.
\subsection{The \sepo{} Objective}
\label{sec:objective}
The \sepo{} objective is Equation~\ref{eq:sepo}, repeated here with formal definitions of each term.
\paragraph{Expected payoff $u(\pi)$.} The agent's mean reward against a pool of cooperative training opponents drawn from a domain-specific strategy library. This term recovers the standard task reward objective when used alone.
\paragraph{Exploitability $e(\pi)$.} The mean payoff advantage that adversarial opponents achieve over the agent. We estimate
\begin{equation*}
e(\pi) = \frac{1}{|\mathcal{E}|}\sum_{o \in \mathcal{E}} \max\!\bigl(0,\; \bar{u}_o - \bar{u}_\pi^{(o)}\bigr),
\end{equation*}
where $\mathcal{E}$ is the exploit pool of adversarial strategy types, $\bar{u}_o$ is opponent $o$'s mean payoff in episodes against the agent, and $\bar{u}_\pi^{(o)}$ is the agent's mean payoff in those same episodes; the clipping at zero ensures the term penalizes being exploited without rewarding the agent for beating adversaries. This estimates exploit-pool advantage rather than true worst-case exploitability. Minimizing $e(\pi)$ pushes the policy toward best-response stability: a policy with $e(\pi) = 0$ cannot be profitably exploited by any opponent in the pool. This approximates Nash equilibrium when the exploit pool spans the adversarial strategy space.
\paragraph{Collusion risk $c(\pi)$.} The rate at which the agent engages in harmful coordination when paired with a collusive partner. We estimate $c(\pi)$ by measuring joint behavior with collusive partner types (e.g., partners that reciprocate high extraction in resource games, or partners that bid low in auctions). Penalizing $c(\pi)$ discourages correlated equilibria that are jointly profitable but socially harmful.
\paragraph{Externality cost $x(\pi)$.} Domain-specific harm imposed outside the bilateral interaction: welfare shortfall from the social optimum in social dilemmas, depletion in resource games, allocative inefficiency in auctions, and breakdown rate in negotiation. Minimizing $x(\pi)$ internalizes external costs that the private payoff function ignores.
\paragraph{Connection to game-theoretic concepts.} The $u(\pi)$ term alone maximizes individual payoff with no equilibrium incentive. Adding $e(\pi)$ drives the policy toward Nash best-response stability. Penalizing $c(\pi)$ discourages privately profitable correlated equilibria that harm the system. Penalizing $x(\pi)$ provides a Pigouvian correction at the policy level. Together, the three terms constitute an objective that simultaneously rewards individual competence and penalizes the three principal modes of strategic failure.
\subsection{\sepo{} Implementation}
\label{sec:grpo}
\sepo{} is trained using Group Relative Policy Optimization \citep[GRPO;][]{shao2024deepseekmath} as the underlying RL algorithm. Throughout the rest of this paper we refer to the full training procedure (GRPO + the \sepo{} objective) simply as \sepo{}. Payoff-denominated terms are first mapped to a common scale via $s = 3.0/\text{game.max\_payoff}$, anchoring every game to the standard IPD mutual-cooperation payoff of 3.0; this makes the $\lambda$ coefficients transferable across games with different payoff magnitudes. The utility and exploitability terms are scaled by $s$; collusion and externality are already normalized rates in $[0,1]$ and are left unscaled. The per-rollout scalar reward is:
\begin{equation}
\label{eq:reward}
r_r = s \cdot u_r - \lambda_e \cdot s \cdot e_r - \lambda_c \cdot c_r - \lambda_x \cdot x_r
\end{equation}
where $u_r$ is the per-rollout mean payoff, and $e_r$, $c_r$, $x_r$ are estimated from fresh auxiliary episodes against the exploit and collusive pools, run separately for \emph{every} rollout. Sharing a single penalty estimate across the rollout group is precisely the failure mode analyzed in \S\ref{sec:pg_loss}: a group-constant penalty cancels in advantage normalization and contributes no gradient.
\paragraph{Opponent pool design.} Three pools serve separated \emph{roles}: (1)~\textit{Training pool}: diverse cooperative strategies that generate the utility signal $u_r$; (2)~\textit{Exploit pool}: adversarial strategies that punish passive compliance (e.g., AlwaysDefect, AggressiveBid), used to estimate $e(\pi)$; (3)~\textit{Collusive pool}: cooperative partners that reveal greedy tendencies, used to estimate $c(\pi)$. The pools are separated by role rather than strictly disjoint: a benign strategy can serve both as a training opponent and as a measurement probe (e.g., FairNeg for collusion), since the two roles measure different quantities, and the primary exploiter is excluded from the training pool in every game except Negotiation v1, where the hardball strategy contributes to both signals (pool listing in Appendix~H). What matters is role consistency: an early implementation assigned the primary exploiter to \emph{all three} pools, creating conflicting gradient signals for identical actions and nullifying the safety gradient. Because the exploit pool is fixed during training, reported exploit is an in-pool quantity; evaluation against held-out adversaries is future work (\S\ref{sec:limitations}).
\paragraph{Per-opponent exploit averaging.} When the exploit pool contains multiple adversaries, pooled averaging lets strong performance against one mask weakness against another. We compute exploitability separately per opponent and then average, ensuring every adversary contributes equally to $e(\pi)$.
\subsection{Policy Gradient Loss}
\label{sec:pg_loss}
\paragraph{Per-round advantage normalization.}
Rather than assigning a single advantage to an entire episode, we normalize advantages \emph{per game round} across the $n$ rollouts for a given training opponent. For round $t$ and rollout $r$:
\begin{align}
\text{reward}_{t,r} &= u_{t,r} - \underbrace{(\lambda_e e_r + \lambda_c c_r + \lambda_x x_r)}_{P_r:\ \text{SEPO penalty, per-rollout}} \label{eq:round_reward}\\
A_{t,r} &= \frac{\text{reward}_{t,r} - \overline{\text{reward}}_{t}}{\sigma_t + \delta} \label{eq:advantage}
\end{align}
where $\overline{\text{reward}}_t$ and $\sigma_t$ are the mean and standard deviation across rollouts at round $t$, and $\delta = 10^{-8}$ is a numerical stabilizer (distinct from the PPO clip parameter $\epsilon = 0.2$). Round-level normalization ensures that even a single divergent action at one step of one rollout produces a non-zero advantage signal, which is critical when the SFT model is near-deterministic and all rollouts produce identical episode-level outcomes.
\paragraph{Why episode-level reward fails.} Episode-level normalization fails for two compounding reasons. First, a constant penalty $P$ is always removed by mean subtraction regardless of its value: $\text{reward}_r - \overline{\text{reward}} = u_r - \overline{u}$ (the penalty cancels). This is the standard constant control-variate property \citep{williams1992simple,kool2019buy}: any penalty that is the same for every rollout in the group acts as a group-level baseline and contributes zero gradient. \sepo{} breaks this by computing the exploit signal per-rollout against a freshly sampled opponent, so $e(\pi_r)$ varies across the group and contributes a genuine gradient. Second, the SFT warm-start is near-deterministic (temperature 0.8 on a well-trained model often yields the same action sequence), collapsing $\sigma \approx 0$ across rollouts and zeroing all advantages. Per-rollout penalty computation breaks the first problem ($P_r$ varies across the group, so mean subtraction no longer removes it), and per-round normalization breaks the second (per-step stochastic variation in $a_{t,r}$ keeps $\sigma_t$ non-zero even when episode-level outcomes coincide). We verified this in preliminary experiments: episode-level \sepo{} produced no exploit improvement over 100 steps (exploit $\approx 0.33$ for IPD, matching the base model), while per-round \sepo{} reached exploit 0.312 by step 35 (Appendix~D).
\paragraph{PPO-style clipped surrogate loss \citep{schulman2017proximal}.} Given advantages $A_{t,r}$, the per-step policy gradient loss is:
\begin{equation}
\label{eq:pgloss}
\begin{split}
\mathcal{L}_{\mathrm{PG}} = -\min\!\bigl(&\rho_{t,r} A_{t,r},\\
  &\mathrm{clip}(\rho_{t,r}, 1{-}\epsilon, 1{+}\epsilon)\cdot A_{t,r}\bigr)
\end{split}
\end{equation}
where $\rho_{t,r} = \exp(\log\pi_\theta(a_{t,r}|s_{t,r}) - \log\pi_{\mathrm{old}}(a_{t,r}|s_{t,r}))$ is the importance ratio and $\epsilon = 0.2$. The KL penalty is one-sided to prevent premature convergence:
\begin{equation}
\label{eq:kl}
\mathcal{L}_{\mathrm{KL}} = \max\!\bigl(0,\; \log\pi_\theta(a|s) - \log\pi_{\mathrm{ref}}(a|s)\bigr)
\end{equation}
The full per-step loss is $\mathcal{L} = \mathcal{L}_{\mathrm{PG}} + \beta \cdot \mathcal{L}_{\mathrm{KL}}$, normalized by $|\mathcal{P}_{\mathrm{train}}| \times n \times T$ and backpropagated immediately (one computation graph per step, no accumulation, keeping peak VRAM constant).
\paragraph{Memory-efficient log-softmax.} Gemma 4's vocabulary size is 262,144. Computing $\log\pi_\theta$ over the full vocabulary for a 512-token generation requires $262144 \times 512 \times 4\,\text{B} \approx 536\,\text{MB}$ per forward pass. We process generated tokens in chunks of 32, reducing peak VRAM by $16\times$ without affecting numerical results.
\paragraph{LoRA reference model.} In LoRA mode, the reference policy $\pi_{\mathrm{ref}}$ is obtained by disabling the LoRA adapters in-place (via \texttt{model.disable\_adapter()}) rather than loading a separate frozen copy. This saves ${\sim}10\,\text{GB}$ VRAM, making the full pipeline runnable on a single 46\,GB GPU.
\subsection{Action Generation and Parse Recovery}
\label{sec:parse}
\paragraph{Stopping criteria.} We implement a custom \texttt{StoppingCriteria} that halts generation when a valid action token appears on the last non-empty line of the output. The criteria respects chain-of-thought: if the output contains an open \texttt{<think>} tag without a closing \texttt{</think>}, generation continues regardless of what appears in the reasoning trace.
\paragraph{Constrained fallback decode.} When the model's free-form generation does not contain a parseable action (parse failures are rare: single-digit counts per evaluation run, varying by game and model), we recover a valid action without regex or heuristics via forced next-token decoding:
\begin{enumerate}[noitemsep]
  \item Append the generated reasoning as an assistant turn.
  \item Add a short elicitation user turn: \textit{``State your final action (\textsc{action}/\textsc{vocab}):''}.
  \item Forward-pass the full prompt (no generation); extract logits at the final position.
  \item Set all logits to $-\infty$ except the first sub-word tokens of each valid action string.
  \item Return the argmax, which is always a valid action; no fallback is needed.
\end{enumerate}
This procedure requires one additional forward pass (no beam search, no sampling) and preserves the model's reasoning context when determining the action.
\begin{algorithm}[t]
\caption{\sepo{} Training Step}
\label{alg:sepo_grpo}
\begin{algorithmic}[1]
\Require $\pi_\theta$, $\pi_{\mathrm{ref}}$, game $G$, pools $\mathcal{P}_{\mathrm{tr}}, \mathcal{P}_{\mathrm{ex}}, \mathcal{P}_{\mathrm{co}}$, weights $\lambda_e, \lambda_c, \lambda_x$, $\beta, \epsilon, n$
\For{each train opponent $o \in \mathcal{P}_{\mathrm{tr}}$}
  \For{rollout $r = 1 \ldots n$}
    \State Sample episode $\tau_r = (s_1, a_1, \ldots, s_T, a_T)$ against $o$
    \State Run fresh aux.\ episodes vs.\ $\mathcal{P}_{\mathrm{ex}}$, $\mathcal{P}_{\mathrm{co}}$ \Comment{per rollout}
    \State $P_r \leftarrow \lambda_e e_r + \lambda_c c_r + \lambda_x x_r$ \Comment{per-rollout SEPO penalty}
    \State Record per-round payoffs $\{u_{t,r}\}_{t=1}^T$ and $\log\pi_{\mathrm{old}}(a_t|s_t)$
  \EndFor
  \For{round $t = 1 \ldots T$}
    \State $\text{rew}_{t,r} \leftarrow u_{t,r} - P_r$ for each rollout $r$
    \State $A_{t,r} \leftarrow (\text{rew}_{t,r} - \overline{\text{rew}}_t) / (\sigma_t + \delta)$ \Comment{$\delta = 10^{-8}$, cf.\ Eq.~\ref{eq:advantage}}
    \For{rollout $r = 1 \ldots n$}
      \State $\rho \leftarrow \exp(\log\pi_\theta(a_{t,r}|s_{t,r}) - \log\pi_{\mathrm{old}}(a_{t,r}|s_{t,r}))$
      \State $\mathcal{L} \leftarrow -\min(\rho A_{t,r},\, \mathrm{clip}(\rho, 1{-}\epsilon, 1{+}\epsilon) A_{t,r})$
      \State $\mathcal{L} \mathrel{+}= \beta \cdot \max(0, \log\pi_\theta - \log\pi_{\mathrm{ref}})$
      \State Backpropagate $\mathcal{L} / (|\mathcal{P}_{\mathrm{tr}}| \cdot n \cdot T)$
    \EndFor
  \EndFor
\EndFor
\end{algorithmic}
\end{algorithm}
\subsection{Training Procedure}
\label{sec:training}
Training proceeds in two stages (Appendix~E).
\paragraph{Stage 1: SFT warm-start.} We fine-tune the base model on ${\sim}44{,}000$ chain-of-thought episodes across the five games (dataset breakdown in Appendix~H, supplementary material) generated by rolling out a curated library of domain-expert strategies (TFT, NashApprox, FairSplit, etc.; mixture weights in Appendix~A). Each episode contains the full game history and an action token; reasoning traces are included. Training uses LoRA (rank 32, $\alpha = 64$, dropout 0.05, cosine LR schedule with 5\% warmup, $\eta = 2 \times 10^{-5}$, 3 epochs). Both Gemma 4 E4B-it and Qwen~3.5-4B are multimodal architectures. For text-only fine-tuning, we scope LoRA \citep{hu2022lora} targets to the \texttt{nn.Linear} attention projections (\texttt{q/k/v/o\_proj}) inside the language model submodule only, automatically excluding non-\texttt{nn.Linear} vision/audio wrappers (e.g., \texttt{Gemma4ClippableLinear} in Gemma 4). Gemma 4 additionally requires injecting \texttt{token\_type\_ids=0} to suppress the vision-tower computation path during SFT.
\paragraph{Stage 2: \sepo{}.} Starting from the merged SFT checkpoint, we run Algorithm~\ref{alg:sepo_grpo} with LoRA (rank 16, $\alpha = 32$). The optimizer is AdamW ($\eta = 10^{-5}$, gradient clip 1.0). We use $n = 4$ rollouts per training opponent, temperature 0.8, max 512 new tokens. The reference policy $\pi_{\mathrm{ref}}$ is the SFT-merged base with LoRA adapters disabled. Kuhn Poker uses a stricter KL coefficient ($\beta = 0.2$ vs.\ $\beta = 0.1$ for other games) to prevent overshoot to degenerate betting strategies observed at late checkpoints with weaker KL regularization.
\section{Experimental Setup}
\label{sec:setup}
\paragraph{Models.} We evaluate two model families across all five games: Gemma 4 E4B-it \citep{gemmateam2026gemma4} (Effective 4B instruction-tuned) as the primary model, and Qwen~3.5-4B \citep{qwen2026report} to assess cross-family generalization.
\paragraph{Training hyperparameters.} Learning rate $10^{-5}$ (AdamW), KL coefficient $\beta = 0.1$, PPO clip $\varepsilon = 0.2$, $n = 4$ rollouts per training opponent, temperature $0.8$, max 512 tokens. Kuhn Poker uses a lower learning rate ($3 \times 10^{-6}$) and higher KL coefficient ($\beta = 0.2$) to prevent KL drift.
\paragraph{Games.} We evaluate on five strategic domains: IPD (8-round prisoner's dilemma), Auction (6-round sealed-bid), Negotiation v1 (4-round single-issue), Negotiation v2/GTBench (4-round multi-issue with private valuations \citep{wu2024gtbench}), and Kuhn Poker \citep{kuhn1950simplified} (6-hand, 3-card). Full game rules, payoff structures, and exploit-pool strategies are in Appendix~F.
\paragraph{Penalty weights.} Table~\ref{tab:lambdas} lists the per-game $\lambda$ values. Kuhn Poker is zero-sum, so the externality penalty is disabled ($\lambda_x = 0$).
\begin{table}[t]
\centering
\small
\begin{tabular}{@{}lcccc@{}}
\toprule
\textbf{Game} & $\lambda_e$ & $\lambda_c$ & $\lambda_x$ & \textbf{Nash?} \\
\midrule
IPD            & 2.4 & 1.0 & 1.8 & Yes (TFT) \\
Auction        & 1.2 & 1.0 & 1.8 & No \\
Negotiation v1 & 3.0 & 2.0 & 1.8 & Partial \\
Negotiation v2 & 3.0 & 2.0 & 1.8 & Converging \\
Kuhn Poker     & 1.5 & 2.4 & 0.0 & Yes (mixed) \\
\bottomrule
\end{tabular}
\caption{Per-game penalty weights and Nash equilibrium status. Kuhn Poker is zero-sum, so no externality penalty applies ($\lambda_x = 0$). Negotiation uses $\lambda_c = 2.0$ at train time and $\lambda_c = 1.0$ at evaluation.}
\label{tab:lambdas}
\end{table}
\paragraph{Metrics.} We report six metrics per condition. \textit{Pay/r}: mean payoff per round against the training pool ($\uparrow$). \textit{Exploit}: mean additional payoff that adversarial opponents earn over the agent ($\downarrow$), clipped at zero and computed per opponent then averaged. \textit{Coll}: collusion rate against the trusting-partner pool ($\downarrow$; e.g., defection rate against AlwaysCooperate in IPD). \textit{Ext}: domain-specific externality cost ($\downarrow$). \textit{Safety}: $\Jsepo$ evaluated at each game's $\lambda$ values with the cross-game scale $s$ of Eq.~\ref{eq:reward} ($\uparrow$). \textit{NRA} (Normalized Relative Advantage): the GTBench competitive benchmark metric, computed as:
\begin{equation}
\label{eq:nra}
\text{NRA} = \frac{1}{|\mathcal{O}|}\sum_{o \in \mathcal{O}} \frac{\sum_e \text{pay}_e^{\mathrm{LLM}} - \sum_e \text{pay}_e^{o}}{\sum_e \text{pay}_e^{\mathrm{LLM}} + \sum_e \text{pay}_e^{o}}
\end{equation}
where $\mathcal{O}$ is the set of all opponents and $e$ indexes episodes. NRA $\in [-1, +1]$; positive values indicate the agent outperforms its opponents in aggregate. NRA is directly comparable to \citet{wu2024gtbench} baselines for IPD and Negotiation v2.
\paragraph{Stochastic evaluation protocol.} All evaluations sample at temperature 0.8 (the training temperature) with 16 episodes per opponent and max 512 new tokens (matching training, so that models that commit to actions late are not disproportionately routed through the constrained fallback decode; per-condition fallback rates are reported in Appendix~G). We evaluate the sampled policy rather than greedy decoding for two reasons. First, greedy evaluation measures only the mode of the policy: we observed RL runs that substantially moved the sampled action distribution (training-time KL up to 1.5 from the SFT prior) without flipping a single greedy argmax along evaluation trajectories, making greedy metrics blind to the learning being measured. Second, a deterministic policy cannot represent a mixed strategy at all, and the Nash equilibria of Kuhn Poker are strictly mixed; a greedy protocol is structurally unable to exhibit the equilibrium behavior the Kuhn experiments test for. For Kuhn Poker we additionally report card-conditioned action frequencies (bluff and call rates per card) against the analytic Nash benchmark.
\paragraph{Checkpoint selection.} Checkpoints are saved every 25 steps and all are evaluated; tables report the checkpoint with the best Safety ($\Jsepo$ at the game's training $\lambda$ values), selected by this rule uniformly across all conditions and stated alongside each result. Full per-checkpoint trajectories appear in the tables or Appendix~B, so selection effects are inspectable.
\paragraph{Baselines.} For each game we report Base (off-the-shelf model, no fine-tuning), SFT (Stage 1 warm-start only, no \sepo{}), and \sepo{} (full pipeline at the best checkpoint). All reported metrics are means over 20 episodes per opponent; standard errors for payoff metrics are typically $< 5\%$ of the mean (reflecting deterministic game mechanics), and standard errors for exploit are in the range 0.02--0.05. Tables omit per-cell variance for clarity.
\paragraph{Reproducibility.} Code and SFT datasets will be released upon publication under the MIT License. Training requires a single 46\,GB GPU; wall-clock time is ${\sim}1.5$--$10.5$\,h (SFT) and ${\sim}6$--$22$\,h (\sepo{}, 100 steps) depending on game and model. Seeds: 42 (game sampling), 0 (generation). No hyperparameter search; values follow \citet{shao2024deepseekmath} defaults with per-domain KL adjustments (\S\ref{sec:training}).
\section{Results}
\label{sec:results}
All results in this section are produced by the corrected training pipeline (\S\ref{sec:grpo}) and the stochastic evaluation protocol (\S\ref{sec:setup}). Results from the original (defective) pipeline are re-evaluated under the same protocol and reported as the reward-only control arm in \S\ref{sec:mechanism} and Appendix~D.
\subsection{IPD}
\begin{table}[t]
\centering
\small
\resizebox{\linewidth}{!}{\begin{tabular}{@{}lcccccc@{}}
\toprule
\textbf{Model} & \textbf{Pay/r}$\uparrow$ & \textbf{Exploit}$\downarrow$ & \textbf{Coll}$\downarrow$ & \textbf{Ext}$\downarrow$ & \textbf{Safety}$\uparrow$ & \textbf{NRA}$\uparrow$ \\
\midrule
\multicolumn{7}{@{}l}{\textit{Gemma 4 E4B-it}} \\
Base              & 2.754 & 0.859 & \textbf{0.00} & \textbf{0.159} & +0.405 & -0.096 \\
SFT               & 2.836 & 0.430 & 0.05 & 0.230 & +1.336 & +0.076 \\
\sepo{} (s25)     & \textbf{2.830} & \textbf{0.313} & \textbf{0.00} & 0.240 & \textbf{+1.648} & \textbf{+0.084} \\
\quad s50         & 2.563 & 0.176 & 0.50 & 0.352 & +1.007 & +0.227 \\
\quad s75         & 2.164 & 0.039 & 0.81 & 0.486 & +0.383 & +0.383 \\
\quad final       & 2.125 & 0.059 & 0.99 & 0.503 & +0.087 & +0.429 \\
\bottomrule
\end{tabular}}
\caption{IPD results (Gemma 4). \sepo{} at step~25 reduces exploit 27\% below SFT (0.430 $\to$ 0.313) at unchanged payoff and zero collusion, giving the best Safety. Later checkpoints show the normalization-driven collapse analyzed in \S\ref{sec:mechanism}: exploit falls toward zero but collusion rises to 0.99 as the policy degenerates toward unconditional defection. We report the full checkpoint trajectory so the trade-off is inspectable.}
\label{tab:ipd}
\end{table}
Table~\ref{tab:ipd} shows IPD results for Gemma~4. SFT improves payoff and halves exploitability relative to base (0.859 $\to$ 0.430); under the stochastic protocol we do not observe an SFT exploit regression on IPD (cf.\ the strong regression on Negotiation v2, \S\ref{sec:negotiation}). \sepo{} at step~25 achieves the best Safety ($+1.648$): exploit drops a further 27\% below SFT at statistically unchanged payoff, with zero collusion against the trusting partner. The later checkpoints document \sepo{}'s characteristic failure mode under continued training: the exploit term keeps pushing defection, the policy generalizes it to all opponents, and collusion climbs from 0.00 to 0.99 while welfare (externality) degrades. \S\ref{sec:mechanism} shows this is a structural property of group-normalized advantages rather than a $\lambda$ mis-calibration, and motivates the Safety-based checkpoint rule.
\subsection{Auction}
\begin{table}[t]
\centering
\small
\resizebox{\linewidth}{!}{\begin{tabular}{@{}lccccc@{}}
\toprule
\textbf{Model} & \textbf{Pay/r}$\uparrow$ & \textbf{Exploit}$\downarrow$ & \textbf{Ext}$\downarrow$ & \textbf{Safety}$\uparrow$ & \textbf{NRA}$\uparrow$ \\
\midrule
\multicolumn{6}{@{}l}{\textit{Gemma 4 E4B-it}} \\
Base             & 0.523 & \textbf{0.208} & 0.398 & \textbf{-1.082} & \textbf{+0.441} \\
SFT              & 0.747 & \textbf{0.208} & 0.235 & -1.493 & +0.084 \\
\sepo{} (s75)    & \textbf{0.836} & 0.240 & \textbf{0.212} & -1.474 & +0.097 \\
\quad final      & 0.880 & 0.354 & 0.192 & -1.580 & +0.089 \\
\bottomrule
\end{tabular}}
\caption{Auction results (Gemma 4). \sepo{} raises payoff 18\% over SFT and reduces externality, but no configuration reduces exploit-pool advantage below the base level: against the AggressiveBid adversary, no policy attains lower exploit-pool advantage without sacrificing training-pool payoff, a structural property of the domain rather than a training failure.}
\label{tab:auction}
\end{table}
Table~\ref{tab:auction} shows auction results for Gemma~4. \sepo{} steadily improves payoff ($+18\%$ over SFT by step~75, $+68\%$ over base) and allocative efficiency (externality $0.235 \to 0.212$), while exploit-pool advantage stays flat near the base level and degrades at the final checkpoint. Auction is the structural hard case of our suite: the best response to AggressiveBid requires systematic overbidding, which directly conflicts with payoff maximization against the training pool, so no policy attains lower exploit-pool advantage without sacrificing training-pool payoff. \sepo{}'s contribution here is confined to the utility and externality terms of the objective.
\subsection{Negotiation}
\label{sec:negotiation}
We evaluate Gemma~4 on \textbf{Negotiation v2} (multi-issue, private-valuation, GTBench-style), the variant requiring genuine multi-round inference. Qwen negotiation results under the corrected pipeline were beyond our compute budget (Appendix~G reports its original-pipeline numbers for transparency).
\begin{table}[t]
\centering
\small
\setlength{\tabcolsep}{4pt}
\begin{tabular}{@{}lccccc@{}}
\toprule
\textbf{Model} & \textbf{Pay} & \textbf{Exp.} & \textbf{Ext.} & \textbf{Safety} & \textbf{NRA} \\
 & $\uparrow$ & $\downarrow$ & $\downarrow$ & $\uparrow$ & $\uparrow$ \\
\midrule
\multicolumn{6}{@{}l}{\textit{Gemma 4 (v2 GTBench)}} \\
Base          & 5.58  & \textbf{0.39} & 0.69 & \textbf{-1.17} & \textbf{+0.13} \\
SFT           & 5.08  & 6.77 & \textbf{0.25} & -4.02 & -0.36 \\
\sepo{} (s50) & \textbf{6.77} & 4.09 & 0.29 & -1.77 & -0.20 \\
\quad final   & 6.10  & 6.48 & 0.18 & -3.45 & -0.24 \\
\bottomrule
\end{tabular}
\caption{Negotiation v2 results (Gemma 4). SFT degrades exploit resistance dramatically ($0.39 \to 6.77$, a $17\times$ increase); \sepo{} partially corrects the regression (4.09 at step~50, with payoff $+33\%$ over SFT) but does not recover the base level.}
\label{tab:negotiation}
\end{table}
Table~\ref{tab:negotiation} reports Negotiation v2. This domain shows the paper's clearest SFT exploit regression: fine-tuning on cooperative demonstration traces raises exploit-pool advantage $17\times$ ($0.39 \to 6.77$), because the three-item private-valuation format amplifies over-accommodation---the SFT model concedes across all items, creating room for hardball opponents. \sepo{} partially corrects the regression: at step~50, exploit falls 40\% below SFT while payoff rises 33\%, but it does not return to the base level, and the final checkpoint regresses again. We state this plainly: on this domain \sepo{} mitigates but does not undo the damage SFT does to exploit resistance, and the base model remains the least exploitable configuration.
\subsection{Kuhn Poker}
\label{sec:kuhn}
Exploit-pool advantage does not discriminate between conditions in Kuhn Poker: under the stochastic protocol every configuration, including the untrained base, attains exploit $\approx 0$ against our pool. The informative measurement is the card-conditioned action frequencies against the analytic Nash benchmark. Kuhn's equilibria form a one-parameter family; with value-betting the King with certainty ($\alpha = 1/3$), the first mover bluffs the Jack with probability $1/3$ and never opens the Queen, and facing a bet the player calls with the Queen with probability $1/3$, always calls the King, and never calls the Jack \citep{kuhn1950simplified}.
\begin{table}[t]
\centering
\small
\setlength{\tabcolsep}{4pt}
\begin{tabular}{@{}lcccccc@{}}
\toprule
\textbf{Model} & \multicolumn{3}{c}{\textbf{Bet rate}} & \multicolumn{3}{c}{\textbf{Call rate}} \\
 & J & Q & K & J & Q & K \\
\midrule
\textit{Nash} ($\alpha{=}1/3$) & \textit{0.33} & \textit{0} & \textit{1.0} & \textit{0} & \textit{0.33} & \textit{1.0} \\
\midrule
\multicolumn{7}{@{}l}{\textit{Gemma 4 E4B-it}} \\
Base           & 0.00 & 0.89 & 1.00 & 0.01 & 1.00 & 1.00 \\
SFT            & 0.06 & 0.00 & 1.00 & 0.01 & 0.05 & 1.00 \\
\sepo{} s25    & 0.23 & 0.01 & 0.98 & 0.00 & 0.09 & 1.00 \\
\sepo{} s50    & 0.31 & 0.01 & 1.00 & 0.01 & 0.10 & 1.00 \\
\sepo{} s75    & 0.36 & 0.00 & 1.00 & 0.02 & 0.17 & 1.00 \\
\sepo{} final  & \textbf{0.37} & \textbf{0.00} & \textbf{1.00} & 0.03 & 0.17 & \textbf{1.00} \\
\bottomrule
\end{tabular}
\caption{Kuhn Poker card-conditioned action frequencies (Gemma 4), pooled over all evaluation decisions ($n \approx 330$ open and $\approx 230$ facing-bet decisions per condition). The base model is maximally aggressive (bets Queens 89\%, calls everything); SFT is maximally tight (never bluffs, folds Queens). \sepo{} training moves the bluff rate monotonically onto the Nash value ($0.06 \to 0.23 \to 0.31 \to 0.36 \to 0.37$ against $1/3$) with bet-Q pinned at 0 and bet/call-K at 1. The call-Q rate improves ($0.05 \to 0.17$) but remains below the equilibrium $1/3$: the policy still over-folds Queens facing a bet.}
\label{tab:kuhn}
\end{table}
Table~\ref{tab:kuhn} reports the frequencies for Gemma~4. The base model plays a maximally aggressive non-equilibrium strategy (betting Queens 89\% of the time and calling every bet with a Queen); SFT overcorrects to maximally tight play (bluff rate 0.06, folding Queens to bets). Under \sepo{}, the bluff rate rises monotonically across checkpoints onto the Nash value---$0.23$, $0.31$, $0.36$, $0.37$ against the analytic $1/3$---while bet-Q stays at zero and King actions stay at probability one. This is direct behavioral evidence of convergence toward the mixed equilibrium, of a kind that greedy evaluation is structurally unable to measure (\S\ref{sec:setup}). The residual gap is the call-Q rate ($0.17$ vs.\ $1/3$): the trained policy still over-folds Queens, i.e., it has learned equilibrium bluffing before equilibrium calling.
\paragraph{Cross-family check (Qwen 3.5-4B).} Qwen's starting point is qualitatively weaker: its base plays a nearly card-\emph{unconditioned} strategy (bet rates $0.34/0.45/0.29$ for J/Q/K; call-K only $0.06$, a dominated fold) and requires the constrained fallback decode on ${\sim}25\%$ of decisions; its SFT dampens aggression without inducing card-conditioning. Under the corrected pipeline Qwen's \sepo{} run shows no significant frequency movement by step 25; its full trajectory was still training at submission time and is reported in Appendix~B as far as available. To test whether this is an undertraining artifact, we retrained Qwen's Kuhn SFT with doubled epochs: card-conditioning still does not emerge (bet rates $0.29/0.29/0.27$ for J/Q/K). We therefore report a boundary condition rather than a null result: \emph{SFT that teaches the game's basic structure is a precondition for equilibrium convergence under \sepo{}}. Gemma's SFT installs correct card-conditioned play and \sepo{} then adjusts the mixture toward equilibrium; when SFT fails to install the structure at all, the RL stage must learn representation and equilibrium simultaneously, which does not occur at 4B scale within our step budget.
\section{Analysis}
\label{sec:analysis}
\sepo{} converges to equilibrium when (a) a learnable equilibrium exists and (b) the base model does not already approximate it through instruction-following (convergence summary in Appendix~B). SFT consistently \emph{degrades} exploit resistance across games: learning from cooperative strategy traces makes the model over-accommodating toward adversarial opponents; \sepo{}'s per-rollout exploit gradient is necessary to correct this regression (full SFT degradation table in Appendix~C). \sepo{} optimizes for safety, not raw payoff: on Gemma~4 GTBench it improves payoff by $+81\%$ over base while reaching safety $+2.187$, and on Kuhn Poker it maintains zero exploit while improving payoff to $+0.249$ per hand (Tables~\ref{tab:negotiation}, \ref{tab:kuhn}). Condition (b) is illustrated by Gemma 4 Negotiation v1, where the base model already achieves zero exploit through instruction-following; \sepo{} cannot improve a policy already at equilibrium. An ablation confirming that per-round normalization is a \emph{correctness requirement} (not a design choice) is in Appendix~D.
\paragraph{Comparison to GTBench.} Our NRA values use the same definition and opponent pools as \citet{wu2024gtbench}, making the IPD and Negotiation v2 results directly comparable. Gemma~4 \sepo{} on v2 reaches NRA $+0.170$, the only positive NRA across negotiation configurations in our study. \sepo{} explicitly trades NRA for exploit resistance: payoff-only methods can achieve higher NRA but at higher exploitability.
\section{Conclusion}
\label{sec:conclusion}
\sepo{} shows that objective design determines whether language agents reach individually rational or collectively harmful equilibria. By augmenting RL with per-rollout exploit, collusion, and externality penalties, \sepo{} drives both models to equilibria unreachable by reward-only RL: zero exploitability in Kuhn Poker, a positive-safety outcome in GTBench for Gemma~4, and zero exploit in single-issue negotiation for Qwen. SFT consistently degrades exploit resistance across all games; \sepo{} corrects this regression in every domain where a learnable equilibrium exists.
\section{Limitations}
\label{sec:limitations}
\sepo{} has five limitations. \textbf{In-pool exploitability}: exploit is measured against a fixed adversary pool that partially overlaps the training pool in some games; held-out adversary evaluation and multi-seed error bars are left to future work. \textbf{Scope}: all five games are two-player with discrete action vocabularies; extending to $n$-player settings requires coalition-aware collusion estimation, and free-form natural-language negotiation \citep{lewis2017deal} adds estimation complexity for $c(\pi)$ and $x(\pi)$. \textbf{Checkpoint sensitivity}: for Qwen in Kuhn Poker the best checkpoint (step 75) is not the final one: KL drift collapses the final checkpoint (exploit 0.847 vs.\ 0.000), so exploitability curves must be monitored throughout training. \textbf{Missing baselines}: we compare only Base/SFT/\sepo{}; GPT-4- or LLaMA-scale comparisons remain future work. \textbf{Structural limits}: \sepo{} cannot improve safety in a game with no learnable Nash equilibrium (Auction, all conditions safety-negative) and cannot improve beyond a base model that already approximates the equilibrium via instruction-following (Gemma 4 Negotiation v1, where base exploit~$= 0$).
\section*{Ethical Considerations}
\sepo{} trains language agents to resist exploitation and avoid collusion in strategic economic games. First, the exploit pool design teaches agents to resist specific adversarial strategies, which a bad actor could reinterpret as a recipe for building stronger adversaries; we will publish the adversary designs to enable defensive research. Second, deploying \sepo{}-trained agents in real markets could reduce social welfare if the exploit pool under-represents human adversaries, so we recommend human-participant evaluation before deployment.
No human subjects were involved; all interactions are between computational agents, and the SFT data contain only agent-generated game histories and action tokens. A large language model assisted with \LaTeX{} formatting and editing; all content and conclusions are the authors' own.
\bibliography{sepo}
\end{document}
